\documentclass{article}
\usepackage{amsmath,amssymb,amsthm,bm}
\usepackage{xcolor}
\usepackage[normalem]{ulem}
\usepackage{hyperref}
\usepackage{enumitem}
\usepackage{booktabs}
\usepackage{listings}
\usepackage{graphicx}
\usepackage{mwe}

\newcommand{\R}{\mathbb{R}}
\newcommand{\pair}[2]{#1 and #2}
\newenvironment{custombox}{\begin{quote}}{\end{quote}}
\newtheorem{theorem}{Theorem}

\title{VisualTeX Document Import Coverage}
\author{VisualTeX}
\date{2026}

\begin{document}
\maketitle

\begin{abstract}
This file is compiled by XeLaTeX and parsed by the VisualTeX Word document importer.
\end{abstract}

\section[Short title]{Common text and inline structures}
Plain text with \textbf{bold}, \textit{italic}, \emph{emphasis},
\underline{underline}, \sout{strike}, \texttt{inline code}, and
\verb|a_b%$| literal content.

A hyperlink is \href{https://example.com}{Example}, a URL is
\url{https://example.com}, and this sentence has a footnote\footnote{A visible note.}.
Colored and transformed content: \textcolor{red}{red text},
\colorbox{yellow}{boxed text}, \fcolorbox{black}{white}{framed text},
\scalebox{1.1}{scaled text}, \resizebox{3cm}{!}{resized text}, and
\rotatebox{2}{rotated text}.

The macros expand to $x\in\R$ and \pair{left}{right}.

\section{Lists, quotations, and code}
\begin{itemize}
  \item First bullet with $a+b$.
  \item Second bullet.
\end{itemize}

\begin{enumerate}[label=(\alph*)]
  \item First numbered item.
  \item Second numbered item.
\end{enumerate}

\begin{description}
  \item[Term] A description-list definition.
\end{description}

\begin{quote}
Quoted text with \textbf{formatting} and inline formula $q=1$.
\end{quote}

\begin{verbatim}
value_1 = 20 % literal comment
\end{verbatim}

\begin{lstlisting}
print("$not_math$")
\end{lstlisting}

\section{Mathematics}
Inline math $E=mc^2$ and \(\alpha+\beta=\gamma\).

\[
\begin{cases}
x^2, & x>0,\\
0, & x\le 0
\end{cases}
\]

\begin{equation}
F=ma
\end{equation}

\begin{align}
a+b &= c,\\
d &= e+f.
\end{align}

\begin{gather}
x=1,\\
y=2.
\end{gather}

\begin{multline}
p+q+r+s+t\\
=u+v+w.
\end{multline}

\begin{flalign}
m&=n&&
\end{flalign}

\begin{alignat}{2}
a&=b &\qquad c&=d
\end{alignat}

\begin{eqnarray}
r&=&s
\end{eqnarray}

\section{Tables, figures, and layout}
\begin{table}
\centering
\begin{tabular}{cc}
\toprule
Name & Value\\
\midrule
Alpha & $\alpha$\\
Beta & $\beta$\\
\bottomrule
\end{tabular}
\caption{A small editable table.}
\end{table}

\begin{figure}
\centering
\includegraphics[width=2cm]{example-image-a}
\caption{A figure placeholder test.}
\end{figure}

\begin{minipage}[t]{0.45\textwidth}
Minipage body with formula $z=3$.
\end{minipage}

\begin{theorem}[Coverage theorem]
Every supported visible structure is preserved or predictably degraded.
\end{theorem}

\begin{custombox}
Custom environment body with $u+v$.
\end{custombox}

\section{External resources and bibliography}
Input fragment text with inline formula $i=1$.

Included fragment text with inline formula $j=2$.


\begin{thebibliography}{9}
\bibitem[Knuth84]{knuth}
Donald Knuth, \emph{The TeXbook}.
\end{thebibliography}

\end{document}
